
\subsubsection{Weight Space Property Prediction}

Unterthiner et al. [2020] established that layer-wise statistics (mean, variance, and higher moments of weight distributions per layer) achieve Spearman $\rho \approx$ 0.9 on simple model zoo tasks. This strong baseline is permutation-agnostic --- it ignores weight space geometric structure entirely, relying only on marginal weight statistics per layer. The Schürholt et al. [2022] model zoo dataset provided a large-scale benchmark of diverse MLP populations, enabling systematic comparison of weight space encoders. Schürholt et al. [2021] proposed Hyper-Representations, self-supervised pretraining on model weight populations, which treat weight populations as datasets for unsupervised learning without directly addressing symmetry structure.

\subsubsection{Equivariant and Invariant Encoders for Weights}

Zhou et al. [2023] introduced Neural Functional Networks (NFN), characterizing the complete class of linear maps equivariant to neuron permutations for MLP weights. Their Neural Functional Transformer (NFT) variant applies attention over weight tokens and achieves permutation equivariance by design. Navon et al. [2023] extended the theoretical analysis to the full symmetry group of MLP weight spaces, including scaling and sign-flip symmetries beyond permutation. For a 2-layer ReLU MLP with hidden dimension h, the scaling symmetry group has dimension h (one free scale per hidden neuron) and the sign-flip symmetry group has order $2^h$. Navon et al.'s DWSNets architecture builds in equivariance to these additional symmetries but is restricted to networks with $M\geq 3$ layers and is incompatible with the M=2 Schürholt MNIST zoo. Critically, while the theoretical characterization is complete, no prior work has empirically measured the geometric diameter of scaling or sign-flip orbits in a real model zoo as cosine distances, nor has any prior work probed whether existing equivariant encoders like NFT are naturally invariant to these orbits in practice.

\subsubsection{Symmetry and Canonicalization in Neural Networks}

Ainsworth et al. [2022] showed that permutation symmetry can be exploited via weight matching to reduce loss barriers between independently trained networks. Entezari et al. [2022] proved that with permutation alignment, loss barriers approach zero under mild conditions. These results motivate orbit-aware processing but focus on permutation, not scaling or sign-flip. In weight space, canonicalization was proposed theoretically by Navon et al. [2023] but not empirically evaluated for property prediction. Kofinas et al. [2024] proposed Universal Neural Functionals extending equivariant processing across architectures, still operating on raw weights without addressing scaling or sign-flip canonicalization.

\textbf{Positioning.} This work adds an empirical measurement layer: orbit characterization establishes which symmetries require canonicalization, how large the orbits are, and whether a given encoder already handles them.

